\section{Conclusion}
This paper studied trigger value selection as an independent design factor in clean-label backdoor poisoning attacks against feature-based malware classifiers.
Regarding attack effectiveness (RQ1), our results show that the values assigned to fixed trigger features strongly affect ASR.
Across EMBER 2018 and EMBER 2024 Win64, aggressive selectors such as MinPopulation and CountAbsSHAP achieved high ASR, particularly under distribution-based distance sampling.
CorrCountAbsSHAP and Q95 also maintained useful attack effectiveness, showing that value selection can change the strength of the backdoor even when the trigger feature set is fixed.

Regarding sanitizer-based detectability (RQ2), our results show that trigger value selection also affects poison recall.
MinPopulation and CountAbsSHAP achieved high ASR, but they were also detected at high rates by Isolation Forest. By contrast, CorrCountAbsSHAP and Q95 reduced poison recall while preserving useful attack effectiveness.
These results show that trigger value selection changes the effectiveness--detectability relationship.
The strongest trigger is not necessarily the most practical if poisoned samples can be removed by sanitization.

Regarding malware-similar benign sampling (RQ3), distribution-based distance sampling generally improved ASR, especially on EMBER 2024 Win64, but it did not remove the effect of value selection. 
Sample selection and value selection are therefore complementary: one determines which benign samples are poisoned, while the other determines the trigger pattern learned by the classifier. 
Overall, trigger value selection is a core design choice that shapes the effectiveness--detectability profile of clean-label malware backdoors.

This study has several limitations. It focuses on feature-space trigger value selection, does not realize triggers in executable PE binaries, and evaluates the proposed selectors only on LightGBM-based classifiers and two PE malware datasets.
Future work should study problem-space realizations, broader models and datasets, stronger end-to-end defenses, and value selectors that jointly optimize attack effectiveness and detectability.